\chapter{Datos ausentes y valores atípicos}\label{cap:ausentes}

\section{Cuánto falta y dónde}

Tras la depuración faltan datos en tres variables (\cref{tab:ausentes}): la cobertura sólo
privada (\res{n-ausente-soloprivado} condados, \res{pct-ausente-soloprivado}\,\%), la incidencia
(\res{n-ausente-incidencia}, \res{pct-ausente-incidencia}\,\%) y el empleo
(\res{n-ausente-empleo}, \res{pct-ausente-empleo}\,\%). Sólo el \res{pct-completos-todo}\,\% de
los condados (\res{n-completos-todo}) está completo en todas las variables, pero la cobertura
sólo privada no se usa como explicativa (es combinación casi lineal de las demás coberturas,
\cref{sec:colinealidad}), de modo que el modelo pierde mucho menos.

\begin{table}[htbp]
\centering
\small
\caption{Datos ausentes tras la depuración y diagnóstico de su mecanismo}\label{tab:ausentes}
\begin{tabularx}{\linewidth}{X r r r l}
\toprule
Variable & $n$ & \% & $p$ ($\chi^2$ estado) & Mecanismo \\
\midrule
Solo seguro privado & \num{609} & \num{20.0} & \num{0.052} & compatible con MCAR \\
Incidencia de cáncer & \num{207} & \num{6.8} & $<\num{0.001}$ & Ausencia estructural \\
Empleo (16+) & \num{152} & \num{5.0} & \num{0.591} & compatible con MCAR \\
\bottomrule
\end{tabularx}
\par\smallskip{\footnotesize\raggedright Contraste MCAR de Little sin la incidencia: $d^2 = \num{62.8}$, \num{74} g.\,l., $p = \num{0.820}$. Con la incidencia: $d^2 = \num{522.3}$, \num{177} g.\,l., $p < \num{0.001}$.\par}
\end{table}


\section{Mecanismo de la ausencia}

El tratamiento adecuado depende del mecanismo \parencite{rubin1987}: si la ausencia es
completamente aleatoria (MCAR), analizar los casos completos no sesga, sólo resta precisión; si
depende de variables observadas (MAR), los casos completos pueden sesgar y conviene imputar con
esas variables; si depende del propio valor no observado (MNAR), ningún método lo resuelve sin
supuestos externos. Se combinaron tres comprobaciones:

\begin{enumerate}
  \item \textbf{Contraste de Little} \parencite{little1988}. Estima por el algoritmo EM la media
  y la covarianza de todas las variables bajo normalidad y compara la media de cada patrón de
  ausencia con la estimada; bajo MCAR, el estadístico
  $d^2=\sum_j n_j(\bar{\mathbf y}_j-\hat{\boldsymbol\mu}_j)^{\top}\hat{\boldsymbol\Sigma}_j^{-1}(\bar{\mathbf y}_j-\hat{\boldsymbol\mu}_j)$
  sigue aproximadamente una $\chi^2$. Con las dos variables de ausencia dispersa no se rechaza
  MCAR ($d^2=\res{little-aleatorio-est}$, \res{little-aleatorio-gl} g.\,l.,
  \res{little-aleatorio-p}); al añadir la incidencia se rechaza con claridad
  ($d^2=\res{little-todo-est}$, \res{little-todo-gl} g.\,l., \res{little-todo-p}).
  \item \textbf{Comparación de medias.} Para cada variable con ausentes, se compara la media de
  las demás entre condados con y sin el dato (Welch, corrección de Holm): ninguna diferencia es
  significativa para el empleo ni para la cobertura sólo privada.
  \item \textbf{Asociación con el estado} ($\chi^2$ de independencia entre la ausencia y el
  estado): la ausencia de la incidencia depende por completo del estado
  (\res{chi-estado-incidencia-p}); la del empleo no (\res{chi-estado-empleo-p}).
\end{enumerate}

\figura{ausentes}{Porcentaje de condados con el dato ausente en los catorce estados con más
ausencias de incidencia: la de la incidencia afecta a estados completos (registros que no
publican datos por condado); las del empleo y la cobertura sólo privada se reparten sin patrón.}

La conclusión es que la incidencia falta por estados enteros (registros que no publican la
cifra por condado): es una ausencia \emph{estructural}, que no depende del valor no observado
(no es MNAR) pero sí del estado, es decir, MAR dado el estado. Las otras dos ausencias son
compatibles con MCAR.

\begin{idea}[Tratamiento elegido.]
El modelo de efectos usa los casos completos de sus variables. Como la ausencia de la incidencia
es MAR, se comprueba que esa elección no altera las conclusiones con una imputación múltiple por
ecuaciones encadenadas \parencite{vanbuuren2011}, con extracciones de la distribución a
posteriori de una regresión bayesiana: \res{m-imputaciones} conjuntos imputados que usan como auxiliares la
respuesta, las demás explicativas y la región, combinados con las reglas de Rubin y los grados de libertad de
\textcite{barnard1999} (\cref{cap:sensibilidad}). La imputación se evalúa sobre los
\res{n-imputacion} condados.
\end{idea}

\section{Valores atípicos}\label{sec:atipicos}

Se buscaron atípicos en tres niveles. Un valor atípico \emph{verificado como real} no se elimina:
eliminarlo sesgaría las estimaciones y subestimaría la variabilidad; lo que importa es si una
observación \emph{influye} desproporcionadamente en el modelo, lo que se evalúa en el
\cref{cap:diagnostico}.

\begin{itemize}
  \item \textbf{Univariantes.} Vallas de Tukey ($Q_1-1{,}5\,\mathrm{RIC}$, $Q_3+1{,}5\,\mathrm{RIC}$)
  y $z$ robusto con la mediana y la MEDA. En la mortalidad hay \res{n-tukey-respuesta} condados
  fuera de las vallas y \res{n-atipicos-respuesta} con $|z|>3{,}5$; todos son condados reales
  (enclaves de altos ingresos con mortalidad muy baja, como \res{min-condado}, o condados rurales
  muy envejecidos o pobres con mortalidad muy alta).
  \item \textbf{Multivariantes.} Distancia de Mahalanobis robusta con el determinante de
  covarianza mínimo \parencite{rousseeuw1999} sobre las \res{mcd-p} explicativas: \res{n-mcd}
  condados superan el cuantil 0,975 de la $\chi^2$ (\res{mcd-corte}), frente a \res{n-mahal-clasica}
  con la distancia clásica. La diferencia ilustra el \emph{efecto máscara}: los atípicos inflan la
  covarianza clásica y se ocultan a sí mismos. Muchos son condados poco poblados con
  composiciones extremas (reservas indígenas, enclaves universitarios), auténticos y no errores.
  \item \textbf{En el modelo}: residuos estudentizados, apalancamiento y distancia de Cook
  (\cref{sec:influencia}).
\end{itemize}
